\src{186}
Considérons alors les accouplements définis par la composition des $\symup{\operatorname{Ext}}$ :
\begin{equation}\tag{6.4}
  \symup{H^{p}(X,F)\times
  \operatorname{Ext}_{\Osh_X}^{n-p}(X;F,\underline{\omega}_X^{n})
  \longrightarrow H^{n}(X,\underline{\omega}_X^{n})}
\end{equation}
et composons-les avec l'homomorphisme $\symup{\gamma}$ de (6.3), on obtient des
homomorphismes fonctoriels, compatibles avec les opérateurs bord :
\begin{equation}\tag{6.5}
  \symup{\operatorname{Ext}_{\Osh_X}^{n-p}(X;F,\underline{\omega}_X^{n})
  \longrightarrow H^{p}(X,F)'}
\end{equation}
(généralisant (5.2)). On vérifie que pour $\symup{p=n}$, on obtient ainsi
l'isomorphisme (6.1). Ceci posé, on a le

\result{THÉORÈME}{3 bis} \uline{Les quatre conditions suivantes sur} $\symup{X}$ \uline{sont
équivalentes pour un entier} $\symup{k\geq0}$ \uline{donné} :
\begin{enumerate}
  \item[\textnormal{i.}] \uline{L'homomorphisme fonctoriel} (6.5) \uline{est un isomorphisme pour}
    $\symup{n-k\leq p\leq n}$.
  \item[\textnormal{ii.}] \uline{On a} $\symup{H^{p}(X,\Osh_X(-m))=0}$ \uline{pour} $\symup{m}$ \uline{grand et pour}
    $\symup{n-k\leq p<n}$.
  \item[\textnormal{iii.}] \uline{Le foncteur} $\symup{H^{p}(X,F)}$ \uline{sur la catégorie des faisceaux
    algébriques cohérents sur} $\symup{X}$ \uline{est coeffaçable pour} $\symup{n-k\leq p<n}$.
  \item[\textnormal{iv.}] \uline{On a}
    $\symup{E^{i}(\Osh_X)=\Extsh_{\Osh_P}^{r-n+i}(\Osh_X,\underline{\omega}_P^{r})=0}$
    \uline{pour} $\symup{0<i\leq k}$.
\end{enumerate}

\noindent DÉMONSTRATION. -- \textnormal{i.}$\symup{\Rightarrow}$\textnormal{ii.} en
vertu du lemme 4, \textnormal{ii.}$\symup{\Rightarrow}$\textnormal{iii.} en vertu du lemme 3,
\textnormal{iii.}$\symup{\Rightarrow}$\textnormal{i.} par un raisonnement standard bien connu,
compte tenu que les deux membres de (6.5) sont alors des foncteurs coeffaçables
pour $\symup{n-k\leq p<n}$ (le premier l'étant en vertu du lemme 4). Enfin prouvons
\textnormal{ii.}$\symup{\Leftrightarrow}$\textnormal{iv.} Cela résulte du corollaire à la
proposition qui suit :

\result{PROPOSITION}{6} \uline{Soit} $\symup{F}$ \uline{un faisceau algébrique cohérent sur} $\symup{X}$, \uline{et
soit} $\symup{i}$ \uline{un entier. Alors pour} $\symup{m}$ \uline{assez grand, on a un isomorphisme} :
\begin{equation}\tag{6.6}
  \symup{H^{i}(X,F(-m))'
  \simeq H^{0}\bigl(X,E^{n-i}(F)(m)\bigr),}
\end{equation}
\uline{où on pose} (comparer n\textsuperscript{o} 3, proposition 5) :
\begin{equation}\tag{6.7}
  \symup{E^{j}(F)=\Extsh_{\Osh_P}^{r-n+j}(F,\underline{\Omega}_P^{r})}
\end{equation}
En effet, en vertu du théorème de dualité sur $\symup{P}$, le premier membre de (6.6)
est isomorphe à
$\symup{\operatorname{Ext}_{\Osh_P}^{r-i}(P;F(-m),\underline{\Omega}_P^{r})}$, donc (6.6) résulte
du lemme 5.

